\documentclass[10pt,a4paper]{amsart}
\usepackage[margin=19mm]{geometry}
\usepackage{amsmath,amssymb,mathtools}
\usepackage[scr=rsfs,cal=euler]{mathalfa}
\usepackage{xfrac}
\usepackage[unicode,hidelinks,bookmarks=false]{hyperref}
\newcommand{\ie}{i.e.\ }
\newcommand{\cf}{cf.\ }
\newcommand{\resp}{respectively\ }
\newcommand{\Subsection}{\subsection}
\providecommand{\nobreakdash}{\nobreak\hbox{-}}
\setlength{\emergencystretch}{2em}
\multlinegap0pt
\usepackage{amssymb}
\usepackage{enumerate}
\usepackage{tikz}
\usepackage{cancel}
\usepackage{mathtools}
\newtheorem{proposition}{Proposition}
\title{Strong persistence index and fluctuations in colon powers of monomial ideals}
\author{Mehrdad Nasernejad, Jonathan Toledo}
\date{Source: arXiv:2604.11475v2 (CC BY 4.0)}
\begin{document}
\maketitle

\noindent\emph{Source and licence.} Strong persistence index and fluctuations in colon powers of monomial ideals. Version arXiv:2604.11475v2 (CC BY 4.0), distributed by arXiv under the Creative Commons Attribution 4.0 International licence, http://creativecommons.org/licenses/by/4.0/, as recorded in the arXiv licence metadata for this exact version and shown on the abstract page https://arxiv.org/abs/2604.11475v2. Full licence notice: \texttt{licenses/arxiv-2604.11475-CC-BY-4.0.txt}.\par\smallskip

\noindent\textit{Editorial scope.} Counted unit: the article's proposition Pro.1 (source0.tex lines 719--724). The article's complete original proof of it is delivered with it (source0.tex lines 726--741, one \texttt{proof} environment of 16 lines). No mathematics is written, repaired or re-derived: the statement and the proof are the article's own lines, in the article's own notation, and no retained line is substituted or re-flowed.  No further source-local context is needed: every cross-reference of every retained line resolves inside the two delivered spans.  Nothing was omitted from any retained span, so the package carries no omission record.  No retained line contains a citation, so the wrapper carries no bibliography and no bibliography entry.  No retained line contains a figure environment, an image-inclusion command or drawing code, so no image is delivered and the package declares no asset dependency.  Editorial formatting only, in addition: an AMS \texttt{amsart} preamble that restates the article's own declarations and macros used by the retained lines, namely the environments \texttt{proposition} on separate counters, together with the standard packages those lines need.  The retained environments print in this document as Proposition 1 in order of appearance, whereas the article prints the same items with the numbering of its own full document.  The complete native source of the article as distributed in the arXiv e-print for this exact version is bundled unchanged as \texttt{sources/198339/source0.tex}.
\begin{proposition} \label{Pro.1}
Let $R$ be a commutative Noetherian ring and let $I$ be an ideal of $R$ whose strong persistence index is $\ell_0$. Then
\[
(I^r : I^s) = I^{r-s} \quad \text{for all } r - s \geq \ell_0.
\]
\end{proposition}

\begin{proof}
Take  $m \in (I^r : I^s)$. Then $m I^s \subseteq I^r$. Take any $h \in I^{s-1}$ and $g \in I$. Since $hg \in I^s$, it follows that
\[
(mh)g = m(hg) \in I^r \quad \text{for all } g \in I.
\]
Thus $mh \in (I^r : I)$ for all $h \in I^{s-1}$. Since $r \geq \ell_0 + 1$, we have
$(I^r : I) = I^{r-1},$
and hence $mh \in I^{r-1}$ for all $h \in I^{s-1}$. This implies that
  $m \in (I^{r-1} : I^{s-1}).$
Repeating this argument inductively $s-1$ times, we obtain
\[
m \in (I^{r-(s-1)} : I^{s-(s-1)}) = (I^{r-s+1} : I).
\]
Since $r - s + 1 \geq \ell_0 + 1$, it follows that $(I^{r-s+1} : I) = I^{r-s}.$ 
Therefore, $m \in I^{r-s}$, and the proof is complete.
\end{proof}

\end{document}
